\section{The tree of function spaces and the one-bit test}
\label{sec:construction}

We first construct questions whose labels are nonempty parts of a joint-evaluation
partition.  Deleting one output direction gives maps with fibers of size at
most two, and a satisfying source assignment supplies labels that satisfy
every such map.  We then specify the distribution of tests, obtaining a
finite weighted game with perfect completeness.  Its soundness will be proved
in the next four sections; the
conversion to the exact alphabets and unweighted format of
Theorem~\ref{thm:main} is deferred to Section~\ref{sec:assembly}.
The shared questions must forget both the names of their output values and
every input coordinate that their evaluations cannot distinguish.  We give
this identification explicitly, since it is responsible for both exact
folding and the later comparison of clause and variable questions.

\subsection{Formal slots and pointwise spaces}

Let the source game be the fixed power of the clause--variable game from
Section~\ref{sec:preliminaries}, and let $t\geq1$ denote its power.
A source left question is an ordered tuple of $t$ clause occurrences.
For a clause occurrence $c$, write $v(c,1),v(c,2),v(c,3)$ for its three
distinct variable IDs in their specified order, and let
\[
 \Omega_c=\mathcal A_c=\{x\in\F^3:x\text{ satisfies }c\}.
\]
These are actual variable values, with the literal signs accounted for in
the satisfaction condition.  Thus $|\Omega_c|=7$, and every coordinate map
$x\mapsto x_j$, and every pair of distinct coordinate maps, is surjective.
A source edge chooses one position in each clause.  Its right question is
the resulting ordered tuple of variable IDs and its answer domain is
$\F^t$.  Coordinates of this domain remain separate even when two IDs agree.

Fix positive integers $d,n_1,\ldots,n_d$, row counts
$\ell_1,\ldots,\ell_d$, and product counts $R_1,\ldots,R_d$.
Let $\mathcal T$ be the rooted ordered tree with a single root at level $d$,
leaves at level $0$, and exactly $n_p$ children at every level-$p$ node.
Write $\mathcal T_p$ for its level-$p$ nodes, and set $K_p=\F^{\ell_p}$.
Every location in this fixed tree is part of its indexing, independently of
the question ID placed there.

Let $E$ be a tuple of source left questions, one at each leaf.  Its elementary slots
are the formal pairs $(C,k)$ with $C\in\mathcal T_0$ and $k\in[t]$.
If the clause occurrence at this slot is $c$, its domain is $\Omega_c$.
All slots are distinct factors of a Cartesian product: repeated clause or
variable IDs impose no equality constraint between their answer coordinates.
For any node $D$, let $\Omega_D(E)$ be the product of the elementary-slot
domains below $D$, and write $\Omega(E)$ for the root domain.

We will also use a mixed tuple $O$ obtained by replacing some leaf left
questions by source right questions along specified source edges.
At a projected elementary slot the domain is $\F$, tagged with the selected
variable ID.  Write $\Omega_D(O)$ and $\Omega(O)$ for the resulting products.
The selected coordinate maps give a surjection
\[
 \pi_D:\Omega_D(E)\longrightarrow\Omega_D(O).
\]
Neither the formal location nor the independence of the answer coordinates
is changed by projection.  For example, the right question $(v,v)$ has the
legal answer $(0,1)$, although a completeness labeling will not use that
answer.

\begin{definition}[Pointwise function spaces]\label{def:spaces}
For any such tuple $Q$ of clause and variable questions, define finite
vector spaces over $\F$ recursively.  At a leaf $C$, set
$H_C(Q)=\F^{\Omega_C(Q)}$, the space of all functions from its answer domain
to $\F$.  For every node $C$, set
\[
 H_C(Q)^2=\Span\{gh:g,h\in H_C(Q)\},
 \qquad (gh)(x)=g(x)h(x).
\]
At a nonleaf node $D$, put
\begin{equation}\label{eq:construction-spaces}
 H_D(Q)=\sum_{C\text{ child of }D} H_C(Q)^2
       \ \subseteq\ \F^{\Omega_D(Q)}.
\end{equation}
Here a function on a descendant domain is extended by ignoring every
other argument.  All sums, products and equalities are equalities of actual
functions on these product domains, rather than formal polynomials.
\end{definition}

\begin{lemma}\label{lem:construction-space-properties}
The spaces in Definition~\ref{def:spaces} contain $1$ and satisfy
$H_C(Q)\subseteq H_C(Q)^2$.  For every proper descendant $C$ of $D$,
\begin{equation}\label{eq:descendant-square-inclusion}
 H_C(Q)^2\subseteq H_D(Q).
\end{equation}
Moreover, pullback by $\pi_D$ injectively identifies $H_D(O)$ with a
subspace of $H_D(E)$ and preserves pointwise multiplication.
\end{lemma}
\begin{proof}
Every elementary domain is nonempty, so extension by ignoring arguments
is injective.  At a leaf, the constant $1$ belongs to the full function
space.  Inductively, any child's constant $1=1\cdot1$ supplies the constant
at its parent.  Multiplication by $1$ gives $H_C(Q)\subseteq H_C(Q)^2$.
For a child, \eqref{eq:descendant-square-inclusion} is the definition.
For a more distant descendant, repeatedly use
$H_C(Q)^2\subseteq H_{D'}(Q)\subseteq H_{D'}(Q)^2$ along the route to $D$.

Each elementary projection onto a selected bit is surjective; their
product $\pi_D$ is therefore surjective.  Consequently pullback is
injective on the entire function space and respects sums and products.
At a leaf its image lies in the full original function space.  Applying
the recursive sum-of-products definition proves the asserted inclusion
at every other node.
\end{proof}

We henceforth use these injective pullbacks without writing their symbols.
No assertion that the summands in \eqref{eq:construction-spaces} form a
direct sum is intended.  In particular their copies of the constant
function coincide.

\subsection{Joint evaluations and canonical questions}

An array at a level-$p$ nonleaf node $D$ is an element
$M_D\in K_p\otimes H_D(E)$, equivalently an ordered list of $\ell_p$ rows
in $H_D(E)$.  At $x\in\Omega(E)$, its evaluation is the vector
$M_D(x)\in K_p$.  The full question function of an array family
$\mathbf M=(M_D)_{D\in\mathcal T,\,D\text{ nonleaf}}$ is
\begin{equation}\label{eq:joint-question-function}
 \Psi_{E,\mathbf M}:\Omega(E)\longrightarrow
       \prod_{p=1}^d\prod_{D\in\mathcal T_p}K_p,
 \qquad x\longmapsto (M_D(x))_D.
\end{equation}
Every nonleaf output occurs in this tuple.

For a nonleaf node $D$ at level $p$ and a nonzero direction $a\in K_p$,
let $Q_{D,a}$ act on the output product in
\eqref{eq:joint-question-function} by the quotient map
$K_p\to K_p/\Span\{a\}$ at the $D$ block and by the identity at every
other block.  Define
\begin{equation}\label{eq:quotient-question-function}
 \Psi^{D,a}_{E,\mathbf M}=Q_{D,a}\Psi_{E,\mathbf M}.
\end{equation}
The quotient merges at most two attained output values.  To turn this
observation into a game constraint, we must first specify which representations
of a joint function give the same question and which answers are legal.

We now turn an array-valued function into a question without retaining
unobservable metadata.  The construction applies to any finite-valued
function $f:\Omega(Q)\to Y$ on a mixed product domain, and in particular
to both functions \eqref{eq:joint-question-function} and
\eqref{eq:quotient-question-function}.  The codomain $Y$ may be a product
of vector spaces and quotient spaces.

An elementary slot $s$ is \emph{unused by $f$} if changing only its
argument never changes $f$.  A clause slot with domain $\Omega_c$ is
\emph{reducible to position $j$} if, for every fixed choice $x_{-s}$ of
all other arguments and every $u,u'\in\Omega_c$ with $u_j=u'_j$,
\begin{equation}\label{eq:global-singleton-factorization}
 f(x_{-s},u)=f(x_{-s},u').
\end{equation}
Thus one fixed position $j$ must work for the whole joint function, for
all choices of the other arguments.

\begin{lemma}\label{lem:canonical-support}
If a used clause slot is reducible to a position, that position is unique.
The following reductions can be performed simultaneously: discard unused
slots; replace each used reducible clause slot by its selected bit; retain
every other used clause slot in full; and retain each used variable slot
as its bit.  They give a surjection
$r_f:\Omega(Q)\to\Omega_f$ and a unique induced function
$\widetilde f:\Omega_f\to Y$ such that $f=\widetilde f\circ r_f$.
The reduced domain and its partition into fibers depend only on the
equality relation
\[
 x\sim_f x'\quad\Longleftrightarrow\quad f(x)=f(x')
\]
on the original product domain, not on names of the output values.
\end{lemma}
\begin{proof}
Suppose the clause slot factors through two distinct positions $j,k$.
Fix all other arguments.  For this section, write the value both as
$g(u_j)$ and as $h(u_k)$.  Pair-surjectivity gives
$g(b)=h(b')$ for every $(b,b')\in\F^2$.  Both functions are constant,
so this section is constant.  This holds for every choice of the other
arguments, making the slot unused.

For each slot let its coordinate reduction be the constant map to a
one-point set, the selected coordinate map, or the identity, according
to the stated rule.  Each is surjective, hence their product $r_f$ is
surjective.  If $r_f(x)=r_f(x')$, replace the slots of $x$ by those of $x'$
one at a time.  At every replacement the value of $f$ is unchanged by
the definition of the reduction at that slot.  This proves that $f$ is
constant on fibers of $r_f$, giving the unique $\widetilde f$.
One-point factors are removed in the notation $\Omega_f$; if all slots
disappear, $\Omega_f$ is a one-point domain.

Both unusedness and \eqref{eq:global-singleton-factorization} are predicates
about pairs of arguments having equal $f$-values.  They are therefore
determined by $\sim_f$.  So are the coordinate reductions, by uniqueness
in the used singleton case.  Surjectivity then shows that the fiber
partition of $\widetilde f$ is determined by the same relation.
\end{proof}

The global quantifier in \eqref{eq:global-singleton-factorization} cannot
be replaced by a sectionwise choice of position.  For example, on
$\Omega_c\times\F$ the function
$f(x,b)=(1+b)x_1+bx_2$ uses one clause position in each fixed-$b$ section,
but has no single position through which it factors globally.  The
canonical reduction retains its clause slot.  Pair-surjectivity proves
uniqueness when a global factorization exists; it does not turn these
different sectionwise choices into one factorization.

The reduced domain has explicit typed coordinates.  At a retained full
clause slot we record its formal location, clause occurrence ID and
triple domain.  At a retained bit slot we record its formal location
and variable ID, with domain $\F$.  When a clause slot is reduced to a
bit, its clause occurrence and selected position are discarded.
Unused locations are discarded.  Retained locations are kept in the fixed
tree-and-slot order, so equal IDs in distinct slots are never merged.
Let
\begin{equation}\label{eq:canonical-partition}
 \mathcal P_f=
 \{\widetilde f^{-1}(y):y\in\im f\}
\end{equation}
be the \emph{unordered} partition of this reduced domain into nonempty
fibers.  The canonical key $\mathfrak k_L(f)$ consists of a left-side
tag, the retained typed coordinates, and $\mathcal P_f$.
Define $\mathfrak k_R(f)$ with a distinct right-side tag.
The codomain $Y$, names $y$ of its values, array coordinates, source
projection positions and any sampling choices are not additional fields
of a key.  The typed coordinates and finite subsets in
\eqref{eq:canonical-partition} have ordinary deterministic finite encodings.

A legal label at this key is a part $P\in\mathcal P_f$.
At a particular representation $f$ of the key, define its restored output
by
\begin{equation}\label{eq:restored-output}
 \operatorname{out}_f(P)=y
 \quad\Longleftrightarrow\quad
 P=\widetilde f^{-1}(y).
\end{equation}
This is a bijection from legal labels to $\im f$.
In particular every restored joint response is realized by one
$x\in\Omega(Q)$: choose a point of the part and lift it through $r_f$.
We never allow separately legal block outputs that do not belong to
the image of the whole joint function.

\begin{lemma}[Sharing across projected domains]\label{lem:canonical-sharing}
Let $O$ be obtained from $E$ by the specified source projections, let
$\pi:\Omega(E)\to\Omega(O)$ be their product, and let
$f:\Omega(O)\to Y$ be any function.  For either side tag,
the keys of $f$ and $f\pi$ coincide literally, with the same parts and
the same restored outputs.
\end{lemma}
\begin{proof}
At an unprojected slot, unusedness and every candidate singleton
factorization are the same for $f$ and $f\pi$, because the projection
of all other arguments is surjective.  At a projected slot, unusedness
is likewise equivalent by surjectivity.  If the slot is used by $f$,
then $f\pi$ factors through its selected position and is used.
Lemma~\ref{lem:canonical-support} makes that position its unique singleton
reduction.  Both representations consequently retain the same formal
location and selected variable ID, with the same bit domain; neither
retains the original clause occurrence or projection position there.
These observations identify their reduced domains.  On that common
domain the two induced functions are identical, since the relevant
product maps are surjective and both lift to $f\pi$.
Their fiber partitions and restored values therefore agree.
\end{proof}

Take as vertices all keys of the indicated side obtainable from mixed
tuples of source questions, arbitrary arrays in their prescribed spaces,
and, on the right, one quotient at a nonleaf node in a nonzero direction.
Every key used by the sampling law below belongs to this set.  Keys never used by a
test have zero incident weight and may be assigned any legal part.
Equivalently, a labeling of the positive-weight vertices can always be
extended to these auxiliary keys.  Because the definition identifies
keys before choosing labels, all such extensions retain the exact sharing
identity of Lemma~\ref{lem:canonical-sharing}.

\subsection{One-bit constraints and perfect completeness}

Fix an array family on $E$, a nonleaf node $D$ at level $p$, and
$a\in K_p\setminus\{0\}$.  Put $f=\Psi_{E,\mathbf M}$ and
$g=Q_{D,a}f$.  The edge joins $\mathfrak k_L(f)$ to
$\mathfrak k_R(g)$ and has the following map on legal parts:
\begin{equation}\label{eq:legal-edge-map}
 \pi_{f,D,a}(P)=
    \operatorname{out}_g^{-1}
       \bigl(Q_{D,a}\operatorname{out}_f(P)\bigr).
\end{equation}
The argument of the inverse belongs to $\im g$, since
$\operatorname{out}_f(P)\in\im f$.
These edge records define the constraints that our game may use.  We specify
their sampling distribution at the end of this section.  Distinct records
may have identical endpoints and may be combined only when their actual
maps agree.

\begin{lemma}\label{lem:construction-two-fibers}
Every map \eqref{eq:legal-edge-map} is a well-defined surjection between
nonempty legal label sets and has at most two preimages per right label.
For a fixed labeling, acceptance of the edge is exactly equality of its
restored outputs through $Q_{D,a}$.
\end{lemma}
\begin{proof}
Restoration is a bijection between parts and attained values on each
side.  Since $g=Q_{D,a}f$, every attained $g$-value has a preimage
among attained $f$-values.  This proves surjectivity and the acceptance
description.  The linear map $Q_{D,a}$ has kernel consisting of the
zero vector and the vector supported at $D$ with value $a$.
Hence every fiber on the entire output space has exactly two elements,
and its intersection with $\im f$ has at most two.
Applying the two restoration bijections proves the claimed fiber bound.
This argument is unchanged if the two question functions retain
different input supports.
\end{proof}

The local label sets have a uniform finite bound once the parameters are
fixed.  Indeed, if $N=|\mathcal T_0|$, every mixed answer domain has size at
most $7^{tN}$, so every legal partition has at most $7^{tN}$ parts.
The canonicalization and each edge map can be computed by enumerating
these finite domains.  Quantitative polynomial-time and weight-removal
details will be given when the parameters have been chosen.

\begin{proposition}[Perfect completeness]\label{prop:completeness}
If the source clause instance has a satisfying assignment, every constraint
just defined is satisfied by one common legal labeling, for every choice
of the positive integer parameters.  Consequently every game supported on
these constraints has value one.
\end{proposition}
\begin{proof}
Let $\xi$ be a satisfying assignment to all source variables.
At a canonical key, assign each retained full clause coordinate the
corresponding satisfying triple from $\xi$, and each retained bit
coordinate its variable value under $\xi$.  These choices form a point
of the reduced domain, even if IDs recur at distinct locations.
Label the vertex by the unique part containing this point.
This prescription depends only on the typed coordinates and partition
in the key, so it gives one consistent label at every shared vertex.

For any original tuple $E$, the assignment $\xi$ gives a legal
point $x_\xi\in\Omega(E)$ whose support reduction is exactly the point
just prescribed at either endpoint.  The restored outputs are therefore
$f(x_\xi)$ and $g(x_\xi)$, respectively.  Their equality through
$Q_{D,a}$ follows from $g=Q_{D,a}f$, so
Lemma~\ref{lem:construction-two-fibers} shows that the edge accepts.
Every constraint therefore accepts, proving the assertion.
\end{proof}

\subsection{Algebraic identities forced by joint answers}

Evaluation tables and folding are standard mechanisms in PCP constructions;
see Bellare, Goldreich and Sudan~\cite[Section~3.3]{BellareGoldreichSudan1998}.
Here the canonical joint partition forces an exact algebraic identity on
every labeling.  It allows an upper array to absorb products of already
displayed descendant rows, including all relations among those products.

\begin{lemma}[Joint folding]\label{lem:folding}
Fix a full array question on $\Omega(Q)$ and a nonleaf node $D$.
For each row index $h$ of $M_D$, choose
\begin{equation}\label{eq:folding-shift}
 b_h=c_h1+
   \sum_{C,u,v}\lambda_{h,C,u,v}
              (M_C)_u(M_C)_v,
 \qquad c_h,\lambda_{h,C,u,v}\in\F,
\end{equation}
where $C$ ranges over any collection of proper nonleaf descendants of $D$
and $u,v$ over their displayed rows.  Let $\widetilde{\mathbf M}$ replace
$(M_D)_h$ by $(M_D)_h+b_h$ and keep every other array fixed.
Then $\widetilde{\mathbf M}$ is a valid array family and its full question
has exactly the same canonical key as that of $\mathbf M$.
For any fixed label at this key, write $y_C$ and $\widetilde y_C$ for
its restored block outputs.  They satisfy
\begin{equation}\label{eq:folding-output}
 \widetilde y_C=y_C\quad(C\neq D),\qquad
 (\widetilde y_D)_h=(y_D)_h+c_h+
   \sum_{C,u,v}\lambda_{h,C,u,v}(y_C)_u(y_C)_v.
\end{equation}
The same statement holds for a question with quotient blocks whenever
the shift uses only unchanged scalar rows actually displayed in that
question; a shift at a quotient block is interpreted through its quotient
map.
\end{lemma}
\begin{proof}
By Lemma~\ref{lem:construction-space-properties}, every term in
\eqref{eq:folding-shift} lies in $H_D(Q)$, proving validity.
On the full output product, define a transformation that fixes every
block other than $D$ and adds the displayed polynomial in those blocks
to block $D$.  Applying this transformation twice is the identity:
its input blocks remain fixed, and addition is in characteristic two.
It is therefore a bijection, and the new joint function is this bijection
composed with the old joint function.  The two functions have identical
equality relations on $\Omega(Q)$.  Lemma~\ref{lem:canonical-support}
gives the same coordinate reduction and the same unordered partition.
Thus a fixed part represents precisely the same set of reduced arguments
in both questions, and evaluating either function on any such argument
gives \eqref{eq:folding-output}.

For quotient blocks use the identical triangular transformation on the
product of displayed quotient spaces.  Its inverse is again itself
because all arguments of the added expression are unchanged displayed
outputs.  The previous proof applies verbatim.  No action on labels is
assumed to be free, and no choice of a representative assignment is
needed to define the identity.
\end{proof}

\begin{lemma}[Joint evaluation of product generators]
\label{lem:joint-evaluation}
Fix a legal response to a full array question and a node $D$.
For any collection $\mathcal C$ of proper nonleaf descendants of $D$, let
\[
 B=\Span\bigl(\{1\}\cup
       \{(M_C)_u(M_C)_v:C\in\mathcal C,\ 1\leq u,v\leq
                                    \ell_{\operatorname{level}(C)}\}\bigr)
       \subseteq H_D(Q).
\]
If $y_C$ are the restored outputs, there is a unique linear functional
$b_y:B\to\F$ such that $b_y(1)=1$ and
$b_y((M_C)_u(M_C)_v)=(y_C)_u(y_C)_v$.
The functional $b_y$ extends to $H_D(Q)^*$.  Every such extension $z$ has
value $1$ on the constant function, and its matrix on these products at each $C$ is
$y_Cy_C^{\mathsf t}$, of rank at most one.
\end{lemma}
\begin{proof}
The restored response has one realizing assignment $x\in\Omega(Q)$.
Point evaluation at $x$ restricts to a linear functional on $B$ and has
the stated generator values.  This proves that every pointwise linear
relation between the generators is respected, including relations
involving different descendants.  The generators span $B$, so this
functional is unique and depends only on the displayed $y_C$.
Extend a basis of $B$ to a basis of $H_D(Q)$ and assign arbitrary values
on the added basis vectors to obtain an extension.  Any extension has
the same generator values, giving the asserted outer product at each descendant.
\end{proof}

The use of all relevant descendant outputs in
Lemma~\ref{lem:joint-evaluation} will allow one upper quotient table to
handle every possible lower node.  A label always supplies those
outputs together; selecting a single descendant is not part of the
vertex key.

\subsection{The distribution of tests}

The low-rank, Hadamard-bucket viewpoint is inspired by Golowich's
construction \cite{Golowich2023}, as mediated by Fei--Minzer--Wang
\cite{FeiMinzerWang2026}.  The sampler here is a recursive-product
redesign with separate independent calls.  We prove its sampling
estimates below; no expansion theorem or identical sampler from those
works is assumed.

We now choose an original tuple $E$, an array family, a node, and an omitted
direction to sample one of the preceding constraints.  Sample a source left
question independently at each leaf, using the source left marginal, and
denote the resulting tuple by $E$.

Independently of $E$, sample a path
$\mathsf P=(D_d,D_{d-1},\ldots,D_0)$ by choosing $D_{p-1}$ uniformly
among the children of $D_p$, successively for $p=d,\ldots,1$.
We next define scalar laws conditional on $E$ and this whole path.
The procedure $\operatorname{Sample}(D_p)$ returns a function in $H_{D_p}(E)$.
At $p=0$ it returns an independent uniform element of $H_{D_0}(E)$.
At $p\geq1$, write $C_*=D_{p-1}$ for the path child and make the following
independent draws:
\[
 U_C\sim\operatorname{Unif}(H_C(E)^2)\quad(C\neq C_*),
 \qquad
 G_h,H_h\sim\operatorname{Sample}(C_*)\quad(1\leq h\leq R_p).
\]
Return
\begin{equation}\label{eq:scalar-sampler}
 \operatorname{Sample}(D_p)
   =\sum_{C\neq C_*}U_C+\sum_{h=1}^{R_p}G_hH_h.
\end{equation}
Its conditional law is denoted $\nu_{D_p}$; the dependence on $E$ and
the path below $D_p$ is suppressed in this notation.
An empty sum over ordinary children is zero.  By
\eqref{eq:construction-spaces}, every returned function lies in the
claimed space.

Every invocation of $\operatorname{Sample}$ uses a new random tape for
the entire recursive call tree.  In particular the two factors in one
product, the factors in different products, and calls made on behalf of
different displayed arrays are independent conditional on $E,\mathsf P$.
Visiting the same geometric node in two calls does not reuse a value.
Uniform functions in ordinary children are drawn directly from the stated
spaces; they do not refer to the displayed arrays in those children.

For every nonleaf path node $D_p$, $1\le p\le d$, independently take a fresh sample
$h_{p,v}\sim\nu_{D_p}$ for each $v\in K_p\setminus\{0\}$ and set
\begin{equation}\label{eq:array-bucket-sampler}
 M_{D_p}=\sum_{v\in K_p\setminus\{0\}}v\otimes h_{p,v}.
\end{equation}
We refer to the summands indexed by $v$ as buckets; $v$ is a bucket
direction in $K_p$, not a row map.
For every nonleaf node off the path, take $M_D$ uniform in
$K_{\operatorname{level}(D)}\otimes H_D(E)$.
All these displayed arrays use separate random tapes.  Thus, conditional
on $E,\mathsf P$, the displayed arrays are independent, including a path
array and an ancestor array whose scalar calls pass through that path
node.  They need not be independent after the common path is averaged out.
The same definitions apply to any mixed tuple $O$ in its own function
spaces.  Later modifications will explicitly replace calls at a specified
cut by uniform draws.

Figure~\ref{fig:tree-cut}, beside the decoder input definitions, depicts the
two cuts and the distinction between a node and separate calls at that node.

Only the canonical key of the joint function
\eqref{eq:joint-question-function} is supplied as the left question.  The
path, buckets, and internal tapes determine its distribution and are not
appended to that key.

To specify all candidate tests on the sampled array family, draw independently
\[
 a_p\sim\operatorname{Unif}(K_p\setminus\{0\}),\qquad 1\leq p\leq d,
\]
independently of the preceding data.  At level $p$ use the quotient
$Q_{D_p,a_p}$; abbreviate it to $Q_{p,a_p}$.  The actual test chooses $p$
uniformly from $[d]$, independently, and uses the edge with left question
$\mathfrak k_L(\Psi_{E,\mathbf M})$, right question
$\mathfrak k_R(\Psi^{D_p,a_p}_{E,\mathbf M})$, and projection
\eqref{eq:legal-edge-map}.  This finite distribution of edge records defines
the preliminary weighted game.  Equivalently, one may draw only the direction
belonging to the chosen level.  Sampling all directions specifies a joint law
of the candidate tests whose individual marginals are precisely this test.

By Proposition~\ref{prop:completeness}, the game has perfect completeness.
The remaining sections prove its soundness by analyzing what a successful
labeling reveals about the sampled arrays.
